%% ─────────────────────────────────────────────────────────────────
%% MLHC 2026 — Research Track Submission
%% Template class: jmlr with pmlr option (PMLR / Proceedings of
%% Machine Learning Research). Download the official style files
%% from: https://www.mlforhc.org/s/MLHC2026-ResearchTrack-Template.zip
%% and place jmlr.cls (and any auxiliary .sty files) in this directory.
%%
%% KEY FORMAT RULES (from MLHC 2026 CfP):
%%   • 11-point Times font  (enforced by jmlr.cls)
%%   • Single column        (do NOT use twocolumn option)
%%   • 10–15 pages          (excluding references and appendix)
%%   • Fully anonymised     (no author names, institution, or links)
%%   • Mandatory section:   "Generalizable Insights" in Introduction
%%   • Submit via OpenReview by April 17, 2026
%% ─────────────────────────────────────────────────────────────────
\documentclass[pmlr]{jmlr}

% ─── Additional packages ─────────────────────────────────────────
% NOTE: jmlr.cls auto-loads: amsmath, amssymb, natbib, graphicx,
%       url, algorithm2e. Only add what is not already loaded.
\usepackage{booktabs}
\usepackage{multirow}
\usepackage{array}
\usepackage{xcolor}
\usepackage{hyperref}
\hypersetup{colorlinks=true, linkcolor=blue, citecolor=blue, urlcolor=blue}

% ─── PMLR metadata (leave volume/year as placeholders for submission)
\jmlrvolume{}
\jmlryear{2026}
\jmlrworkshop{Machine Learning for Healthcare (MLHC) 2026}

% ─── Title (short title appears in header) ───────────────────────
\title[PD Detection from Bilateral IMU via Cross-Attention Transformers]{%
  Parkinson's Disease Detection from Bilateral Wrist-Worn Smartwatch
  IMU Data Using Dual-Channel Cross-Attention Transformers}

% ─── Authors — ANONYMISED FOR SUBMISSION ─────────────────────────
% Replace the block below with your real author block in the
% camera-ready version only. Do NOT de-anonymise before acceptance.
\author{%
  \Name{Anonymous Author(s)} \Email{anonymous@institution.edu}\\
  \addr Anonymous Institution
}

% ─── Editor (leave blank for submission) ─────────────────────────
\editor{}

% ═════════════════════════════════════════════════════════════════
\begin{document}
\maketitle

% ─── Abstract ────────────────────────────────────────────────────
\begin{abstract}
Parkinson's Disease (PD) is a chronic neurodegenerative disorder
characterised by motor symptoms including tremor, bradykinesia,
postural instability, and freezing of gait (FoG). Clinical diagnosis
currently relies on evaluation by a physician or specialist, which is
time-consuming, subject to inter-rater variability, and inaccessible
in many settings. Advanced wearable sensors offer a promising
alternative by enabling continuous, objective monitoring of motor
function in daily life. We propose a Dual-Channel Cross-Attention
Transformer that processes bilateral wrist-worn IMU signals from the
PADS dataset (469 participants: 291 PD, 79 HC, 99 DD) to classify PD
against both Healthy Controls (HC) and Differential Diagnosis (DD)
conditions. Rather than a single three-class objective, we decompose
the problem into two focused binary classification heads, enabling the
model to learn the HC/PD and PD/DD decision boundaries independently.
Patient-level stratified 5-fold cross-validation is applied throughout
to obtain less biased estimates of true out-of-sample performance. The
base model achieves a combined mean accuracy of $92.83\%$ across
folds. We further demonstrate that self-supervised pre-training with
contrastive infoNCE loss reaches $91.02\%$ combined accuracy with only
$20\%$ of labelled data, highlighting the potential of label-efficient
learning for clinical applications. The optimised model achieves a
mean inference time of $48.32$~ms per window on a Raspberry Pi CPU,
confirming viability for real-time edge deployment on smartwatch
hardware. Early and continuous monitoring of PD motor symptoms enabled
by this framework could support timely clinical intervention and reduce
reliance on infrequent specialist visits.
\end{abstract}

% ─── Keywords ────────────────────────────────────────────────────
\begin{keywords}
Parkinson's Disease, IMU, wearable sensing, cross-attention transformer,
self-supervised learning, infoNCE, nested cross-validation, edge deployment
\end{keywords}


% ═════════════════════════════════════════════════════════════════
\section{Introduction}
\label{sec:intro}
% ═════════════════════════════════════════════════════════════════
% size of the introduction should be around <=2 pages we can still add text equivalent to the size of generalized text, we can extend the key contributions,i don't want generalization para, it seems redaundent with key contribution

%these parts are written by me 
Parkinson's is a chronic neurodegenarative disorder that affects 10 million people around the world.
Early symptoms of Parkinson's involve various moto diseases such as  resting tremor, bradykinesia, rigidity,
postural instability, and freezing of gait (FoG). Parkinson's though irreversible constant monitoring and treatment plan can help delay disease progression and increase life expectancy as well as increase the quality of life of a Parkinson's patient.

Currently parkinson's disease diagnosis require direct involvement of health care professionals and clinical rating scale (e.g., UPDRS) which can be inadequate and inherently subjective.

% this para is important, we are not noly doing HC vs PD but PD vs DD also 
Parkinson's shares common motor symptoms with various neurological disorders such as Essential Tremor, Multiple System Atrophy, progressive supranuclear palsy etc, making differentiating
parkinson's disease can be challenging due to it's underlying similarity in motor symptoms, which can often delay the treatment procedure until the severity increases. Distinguishing Parkinson's disease(PD) from differetial diagnosis (DD) due to similar motor patterns remains a clinically a challenging task.

This creates necessity for a system that can detect any abnormalities as well as contrast parkinson's with other differential diagnosis.  %have to make line line longer

The application of advanced wearable sensors in personal devices (e.g. smartwatch,smartphone) has grown rapidly. These wearable sensors can record and monitor physiological and movement data. The smartwatch sensor can record Inertial Measurement Units (IMU) signal which can be  valuable in detecting motor symptoms related to parkinson's disease. However mining meaningful patterns and trends is inherently difficult due to the nature of sensor data,which can often contain noise and signal distortions due to environmental variables.

Parkinson's can show asymmetry in motor symptoms, smartwatches can be worn in both wrist and measure the motor signal without affecting the daily activity of the patient. The goal is to develop a method that can capture the underlying motor asymmetry of parkinson's disease.  %till this point 

Deep learning approaches have shown strong performance in Human Activity Recognition (HAR) and medical time series classification. Sequential models such as LSTMs capture temporal dependencies in
sensor signals, while Transformer architectures with self-attention mechanisms are well suited to modelling long-range temporal patterns. In parallel, self-supervised contrastive representation learning has emerged as a powerful learning paradigm from unlabelled data, which is especially relevant in clinical settings where labelled data is coarse.

This study implements a cross-attention based Transformer framework for Parkinson’s disease detection using bilateral wrist-worn IMU data. The key contributions are: 
\begin{itemize}
  \item A bilateral cross-attention architecture that explicitly models the motor asymmetry between the left and right wrists, which is a clinically established characteristic of PD.
  \item A hierarchical binary classification strategy by decomposing the three-group problem into two focused tasks (HC vs.~PD and PD vs.~DD), resulting in higher accuracy in separating parkinson's from other motor diseases.
  \item A self-supervised pre-training stage using infoNCE contrastive loss that achieves label-efficient transfer, reaching saturation with as few as $20\%$ of available labels.
  \item Deployment of the optimised model on a Raspberry Pi CPU, demonstrating real-time inference at $48.32$~ms per window and validating feasibility for edge deployment on smartwatch hardware.

\end{itemize}



% Parkinson's Disease is a chronic neurodegenerative disorder affecting
% over 10 million people worldwide~\citep{dorsey2018}. It is characterised by
% progressive impairment of the dopaminergic system, leading to hallmark
% motor symptoms including resting tremor, bradykinesia, rigidity,
% postural instability, and freezing of gait (FoG). FoG episodes can
% persist for a few seconds or last several minutes, significantly
% increasing fall risk and reducing quality of life. Currently, PD is
% measured and analysed by a physician or specialist through subjective
% clinical rating scales such as the Unified Parkinson's Disease Rating
% Scale (UPDRS), which are expensive, infrequent, and inherently subject
% to inter-rater variability.

% Advanced wearable sensors, particularly Inertial Measurement Units
% (IMU) embedded in smartwatches, present a compelling opportunity to
% address these challenges. Smartwatch-based systems can capture
% continuous, objective motor data during daily activities without
% disrupting the patient's routine. However, mining meaningful patterns
% and trends from such medical time series data is inherently
% challenging: signals are noisy, high-dimensional, and reflect subtle
% motor irregularities that must be disentangled from normal movement.
% Moreover, because of underlying similarity in motor patterns,
% distinguishing PD from Differential Diagnosis (DD) conditions---other
% neurodegenerative disorders that produce Parkinson's-like
% symptoms---remains a significant clinical and computational challenge.

% Deep learning approaches have shown strong performance in Human
% Activity Recognition (HAR) and medical time series classification.
% Sequential models such as LSTMs capture temporal dependencies in
% sensor signals, while Transformer architectures with self-attention
% mechanisms are well suited to modelling long-range temporal patterns.
% In parallel, self-supervised contrastive representation learning has
% emerged as a powerful paradigm for learning from unlabelled data,
% which is especially relevant in clinical settings where obtaining true
% labels requires expert annotation and is therefore costly.

% We propose a Dual-Channel Cross-Attention Transformer framework for PD
% detection from bilateral wrist-worn IMU data. The key contributions are:
% \begin{itemize}
%   \item A dual-stream cross-attention architecture that explicitly
%         models the motor asymmetry between the left and right wrists,
%         which is a clinically established characteristic of PD.
%   \item A hierarchical binary classification strategy that decomposes
%         the three-group problem into two focused tasks (HC vs.~PD and
%         PD vs.~DD), yielding a $4.48$~pp improvement in combined
%         accuracy over a three-class baseline.
%   \item A self-supervised pre-training stage using infoNCE contrastive
%         loss that achieves label-efficient transfer, reaching
%         competitive performance with as few as $20\%$ of available
%         labels.
%   \item Deployment of the optimised model on a Raspberry Pi CPU,
%         demonstrating real-time inference at $48.32$~ms per window and
%         validating feasibility for edge deployment on smartwatch
%         hardware.
% \end{itemize}

% % ─────────────────────────────────────────────────────────────────
% % MANDATORY FOR MLHC: Generalizable Insights
% % Every MLHC Research Track paper MUST include this subsection
% % in the Introduction. Describe what insights from your work
% % generalise BEYOND this specific task/dataset.
% % ─────────────────────────────────────────────────────────────────
% \subsection*{Generalizable Insights}
% \label{sec:insights}

% This work provides three generalizable insights for the machine
% learning for healthcare community:

% \begin{enumerate}
%   \item \textbf{Bilateral sensing via cross-attention is clinically
%         motivated and broadly applicable.} Many neurodegenerative and
%         musculoskeletal conditions exhibit lateralised motor
%         asymmetry. The dual-stream cross-attention design used here
%         to model contralateral wrist relationships is directly
%         applicable to any clinical wearable sensing scenario where
%         bilateral signals carry differential diagnostic information.

%   \item \textbf{Hierarchical decomposition of a multi-class clinical
%         problem improves performance when class boundaries are
%         qualitatively different.} The $4.48$~pp gain from
%         decomposing HC/PD/DD into two focused binary tasks suggests
%         that this strategy generalises to other clinical settings
%         with heterogeneous class separability, such as staging of
%         diseases or distinguishing comorbid conditions.

%   \item \textbf{Self-supervised pre-training on clinical IMU data
%         substantially reduces label requirements.} With only $20\%$
%         of labelled data, performance exceeds $91\%$, and further
%         labelling yields diminishing returns. This finding
%         generalises to any clinical wearable scenario where expert
%         annotation is scarce---a pervasive bottleneck in medical
%         machine learning.
% \end{enumerate}  

The remainder of this paper is organised as follows.
Section~\ref{sec:related} reviews related work.
Section~\ref{sec:dataset} describes the PADS dataset.
Section~\ref{sec:method} details the proposed methodology.
Section~\ref{sec:experiments} presents experimental results and
ablation studies. Section~\ref{sec:discussion} discusses findings and
limitations. Section~\ref{sec:conclusion} concludes.


% ═════════════════════════════════════════════════════════════════
\section{Related Work}
\label{sec:related}
% ═════════════════════════════════════════════════════════════════

\subsection{IMU-Based Parkinson's Disease Detection}

Wearable IMU sensors have been widely explored for objective PD
assessment. Early work focused on handcrafted features extracted from
accelerometer signals to detect tremor and bradykinesia~\citep{patel2009}.
More recently, deep learning approaches have demonstrated competitive
performance without manual feature engineering. Convolutional Neural
Networks (CNNs) have been applied to IMU time series for gait and
tremor classification, and recurrent architectures such as LSTMs have
been used to capture sequential dependencies in motor signals.
Systematic reviews of deep learning with wearable sensors for
PD~\citep{sigcha2020} highlight the diversity of experimental protocols and
the difficulty of cross-study comparison due to variations in dataset
size, task design, and validation strategy. A key limitation across
this body of work is the predominant use of unilateral sensing,
despite the well-established motor asymmetry of PD, where symptoms
onset and progress more severely on one side of the body. Our work
explicitly exploits bilateral sensing through a cross-attention
mechanism designed around this clinical observation.

\subsection{Transformer Architectures for Medical Time Series}

Transformer models, originally proposed by \citet{vaswani2017} for
natural language processing, have been increasingly applied to time
series classification. Their self-attention mechanism allows the model
to attend to distant timesteps simultaneously, which is advantageous
for capturing the irregular, episodic nature of motor symptoms such as
FoG. Dual-stream or two-stream transformer designs have been proposed
for HAR, using parallel encoders to process different sensor
modalities or spatial streams before fusion~\citep{xiao2022}. However,
existing dual-stream designs typically fuse streams through simple
concatenation or addition, without explicitly modelling the
contralateral relationship between left and right motor signals. Our
architecture fills this gap by replacing independent fusion with
cross-attention between the two streams.

\subsection{Self-Supervised Contrastive Learning for Time Series}

Self-supervised contrastive representation learning has gained
significant traction as a method for learning from unlabelled data.
In the contrastive framework, a model is trained to map an anchor
sample $x_i$ and its augmented positive counterpart $x_i^+$ to nearby
embeddings, while pushing embeddings of negative samples $x_j$
($j \neq i$) further apart. The infoNCE loss~\citep{oord2018} formalises
this objective by treating all other samples in a batch as negatives,
enabling the model to contrast several negative samples simultaneously,
which is more efficient for representation learning than single-pair
objectives. TS-TCC~\citep{eldele2021} proposes a time-series-specific
framework combining temporal and contextual contrasting with strong
and weak augmentations, demonstrating competitive label efficiency on
HAR and medical datasets. These methods show that self-supervised
contrastive representation learning can learn general representations
from unlabelled medical time series that transfer effectively to
downstream tasks with limited annotations. However, these methods have
not been evaluated on bilateral IMU data for PD detection, where the
augmentation strategy must preserve clinically meaningful asymmetric
motor patterns. Our work addresses this gap.

\subsection{Nested Cross-Validation in Clinical Machine Learning}

A persistent problem in medical machine learning is optimistic bias:
when hyperparameter selection and performance estimation are performed
on the same data split, reported metrics overestimate true
out-of-sample performance~\citep{varma2006}. Nested cross-validation
addresses this by separating the model selection loop (inner folds)
from the performance estimation loop (outer folds), providing less
biased estimates of how well predictive performance generalises to new
data points. This bias-variance tradeoff is especially important in
clinical settings with small patient cohorts, where single-fold or
holdout validation can produce highly variable and misleading
estimates. Despite its importance, many published PD detection studies
report performance from a single train/test split or flat k-fold,
making results difficult to compare and potentially overoptimistic.
In this work, we apply Optuna-based hyperparameter search on an inner
validation fold (Fold~1) and evaluate the selected configuration
across all five outer folds using patient-level stratified splitting
to prevent data leakage.


% ═════════════════════════════════════════════════════════════════
\section{Dataset}
\label{sec:dataset}
% ═════════════════════════════════════════════════════════════════

\subsection{Overview}

We use the PADS (Parkinson's Disease Smartwatch) dataset from
PhysioNet~\citep{varghese2024}. It contains sensor data from bilateral
wrist-worn smartwatch recordings of 469 participants performing 10
standardised motor tasks. The dataset comprises three groups: 291
patients with Parkinson's Disease (PD), 79 Healthy Controls (HC), and
99 patients with Differential Diagnosis (DD) conditions that show
Parkinson's-like motor symptoms but are caused by different
neurodegenerative diseases. IMU data is chosen for this task given its
established effectiveness in detecting motor symptoms such as tremor
and bradykinesia.

\subsection{Sensor Modality}

Each participant wears a smartwatch on both wrists, providing six-axis
IMU measurements per wrist. Three axes are from the accelerometer
($a_x, a_y, a_z$), which detects linear movements, and three axes are
from the gyroscope ($\omega_x, \omega_y, \omega_z$), which records
angular movements. Signals are sampled at $100$~Hz.
Figure~\ref{fig:sensor_viz} visualises the sensor data of a
participant performing the \textit{HoldWeight} task, showing distinct
patterns indicative of tremor and bradykinesia.

\begin{figure}[!ht]
  \centering
  \includegraphics[width=0.5\textwidth,height=0.25\textheight,keepaspectratio]{image5.png}
  \caption{IMU sensor data from a participant performing the
           \textit{HoldWeight} task. Distinct temporal patterns in the
           accelerometer and gyroscope channels are indicative of motor
           symptoms.}
  \label{fig:sensor_viz}
\end{figure}

\subsection{Motor Tasks}

Ten standardised motor tasks are recorded from clinical trials:
\textit{CrossArms}, \textit{DrinkGlas}, \textit{Entrainment},
\textit{HoldWeight}, \textit{LiftHold}, \textit{PointFinger},
\textit{Relaxed}, \textit{StretchHold}, \textit{TouchIndex}, and
\textit{TouchNose}. Tasks \textit{Relaxed} and \textit{RelaxedTask}
are 20~seconds in duration; all remaining tasks are 10~seconds.


% ═════════════════════════════════════════════════════════════════
\section{Methodology}
\label{sec:method}
% ═════════════════════════════════════════════════════════════════

\subsection{Problem Definition}

This study aims to develop a method for classifying PD from HC and
from DD using time-series data collected from wearable IMU sensors.
Each input sample contains left ($L$) and right ($R$) wrist signal
pairs, denoted as $\mathbf{X} = (\mathbf{X}^L, \mathbf{X}^R) \in
\mathbb{R}^{T \times 2C}$, where $C{=}6$ denotes the number of sensor
channels per wrist, $2C{=}12$ is the total channel count across both
wrists, and $T{=}256$ is the number of timesteps per sample. Each
sample is associated with a class label $y \in \{\text{HC}, \text{PD},
\text{DD}\}$.

Instead of a single three-class objective, we divide the problem into
two distinct binary classifiers. The model learns a mapping function
$f : \mathbb{R}^{T \times 2C} \rightarrow [0,1]^2$, predicting
probabilities for HC vs.~PD (sample space $y \in \{\text{HC},
\text{PD}\}$) and PD vs.~DD (sample space $y \in \{\text{PD},
\text{DD}\}$). The PD--DD boundary is especially challenging due to
the overlap of motor patterns between the two groups.

\subsection{Data Preprocessing}

\subsubsection{Label Encoding}

A hierarchical labelling strategy supports two independent binary
classification heads. DD samples are masked ($-1$) and ignored during
loss computation for the HC-PD head; HC samples are masked ($-1$) for
the PD-DD head. This ensures each head is trained only on relevant
samples, which also prevents data leakage between heads.

\subsubsection{Patient-Level Stratified K-Fold Cross-Validation}

A patient-level stratified 5-fold cross-validation strategy is
applied. The stratified approach ensures equal class distribution in
each fold. K-fold cross-validation is applied to obtain less biased
estimates of true out-of-sample prediction performance, reducing
variance in evaluation across different data partitions. A
patient-level split is enforced: all windows from a given patient are
assigned entirely to either the training or the validation set,
preventing data leakage across patient boundaries.

\subsubsection{Downsampling}

Raw signals are resampled from $100$~Hz to $64$~Hz via polyphase
resampling. Frequencies above $64$~Hz are irrelevant to HAR
activities, since human motor signals are predominantly below
$20$~Hz, so the reduction in sampling rate incurs no loss of
diagnostic information.

\subsubsection{Band-Pass Filtering}

Human voluntary movements and clinically relevant motor symptoms such
as tremor and FoG are concentrated in the low-frequency,
high-sensitivity band of $0.1$--$20$~Hz. A fourth-order Butterworth
band-pass filter with passband $[0.1,\,20]$~Hz is applied using
zero-phase forward-backward filtering to avoid phase distortion. The
low-frequency cutoff removes DC drift; the high-frequency cutoff
suppresses sensor noise above the motor-relevant band.

\subsubsection{Overlapping Window Segmentation}
\label{sec:windowing}

IMU signals are segmented into consecutive windows of $256$
timesteps ($4$~seconds at $64$~Hz). The participant counts are
disproportionate (HC: 79, PD: 291, DD: 99); uniform non-overlapping
sampling produces severe class imbalance, biasing the model toward PD.
To mitigate this, an overlapping sampling strategy is applied:
minority classes receive higher window overlap to generate more samples
from the same recording duration. We applied $70\%$ overlap to HC
samples, $0\%$ to PD, and $65\%$ to DD. This produced approximately
$8{,}000$ windows per class in the training set ($25{,}640$ total),
distributed evenly across folds via stratified sampling. The number of
windows for a signal of length $L$ with window size $W$ and step size
$s$ is:
\begin{equation}
  N_w = \left\lfloor \frac{L - W}{s} \right\rfloor + 1
\end{equation}

\subsection{Dual-Channel Cross-Attention Transformer}
\label{sec:base_model}

The architecture processes left and right wrist window samples
independently through parallel encoders. The design is motivated by
the well-established motor asymmetry of Parkinson's Disease: symptoms
typically occur more severely on one side, making cross-attention
between left and right wrist signals clinically meaningful.
Figure~\ref{fig:architecture} illustrates the full pipeline.

\begin{figure}[!ht]
  \centering
  \includegraphics[width=0.7\textwidth,height=0.25\textheight,keepaspectratio]{cross_attn.png}
  \caption{Dual-channel cross-attention Transformer. Left and right
           streams are projected, position-encoded, and refined
           through $N$ cross-attention layers before pooling and
           dual-head classification.}
  \label{fig:architecture}
\end{figure}

\subsubsection{Input Projection and Positional Encoding}

Each wrist's $C{=}6$ channel input is linearly projected to a
$d$-dimensional embedding space. Sinusoidal positional encodings are
applied to inject temporal order into the embeddings, following the
standard Transformer methodology~\citep{vaswani2017}.

\subsubsection{Cross-Attention Layers}

Each of the $N$ stacked cross-attention layers has three sequential
sub-layers. In the cross-attention sub-layer, each wrist's embedding
acts as Queries ($Q$) and attends over the opposite wrist's embedding
as Keys ($K$) and Values ($V$), allowing each stream to gather
complementary information from the contralateral wrist. The output is
followed by a residual connection and layer normalisation. For the
left stream, where $\mathbf{H}^L$ queries into $\mathbf{H}^R$:
\begin{equation}
  \mathbf{H}^L \leftarrow \mathrm{LayerNorm}\!\left(
      \mathbf{H}^L + \mathrm{MHA}(\mathbf{H}^L,\, \mathbf{H}^R,\,
      \mathbf{H}^R)
  \right)
\end{equation}
The right stream is updated symmetrically with $\mathbf{H}^R$ as
queries over $\mathbf{H}^L$. A self-attention sub-layer then refines
each stream's own embedding. Finally, a position-wise feed-forward
network applies a two-layer MLP with ReLU activation.

\subsubsection{Pooling, Fusion, and Classification}

After $N$ attention layers, global average pooling is applied to each
stream, yielding $\mathbf{z}^L, \mathbf{z}^R \in \mathbb{R}^d$.
These are concatenated to form $\mathbf{z} = [\mathbf{z}^L \|
\mathbf{z}^R] \in \mathbb{R}^{2d}$, which is fed into two independent
classification heads---one for HC-PD and one for PD-DD---each
implemented as a single linear layer followed by softmax.

\subsection{Training Procedure}

\subsubsection{Objective Function}

The cross-entropy loss is applied in each classification head,
computed only over samples with valid (non-masked) labels for that
head. The joint loss is the average of the two masked losses:
\begin{equation}
\mathcal{L} = \frac{1}{|\mathcal{S}|} \sum_{s \in \mathcal{S}}
\mathcal{L}_s, \qquad
\mathcal{L}_s = -\sum_{c=1}^{2} y_{s,c} \log \hat{p}_{s,c}
\end{equation}

\subsubsection{Optimisation}

The model is trained with AdamW (weight decay $0.01$), gradient
clipping at unit norm, and a ReduceLROnPlateau scheduler (factor
$0.5$, patience $5$ epochs). Training runs for up to $100$ epochs
with batch size $32$.

\subsection{Self-Supervised Pre-training}
\label{sec:ssl}

We apply self-supervised contrastive representation learning as a
pre-training stage before fine-tuning on labelled data. True labels
for clinical samples require annotation by a specialist, making
labelled data scarce and expensive. By first learning general signal
representations from unlabelled windows, the model can transfer
useful structural information to the downstream classification task.

Three types of samples are required in contrastive learning: an anchor
sample $z_i$, a positive sample $z_i^+$ created by augmenting the
anchor, and negative samples $z_j$ ($j \neq i$) drawn from the same
batch. Crucially, only the encoder weights are optimised during
pre-training; the classification heads are not updated and are trained
exclusively during the subsequent fine-tuning stage.

Two augmentation strategies are applied in alternation with $50\%$
probability: temporal warping~\citep{um2017} and Gaussian noise
($\sigma{=}0.05$). Temporal warping applies non-linear transformations
to the time axis via cubic spline interpolation with $k{=}4$ interior
knots displaced from $\mathcal{N}(1.0,\,0.3^2)$, applied
independently per IMU channel with pinned endpoints to avoid boundary
artifacts. Samples are chosen randomly without using any labels.

The infoNCE (Information Noise-Contrastive Estimation)
loss~\citep{oord2018} compels the embedding of anchor $z_i$ to be closer
to its positive $z_i^+$ while pushing it away from all other samples
in the batch ($i \neq j$). Using the batch as a source of negatives
allows the model to contrast several negative samples simultaneously,
which is more efficient for learning representations than single-pair
approaches:
\begin{equation}
\mathcal{L}_{\mathrm{NCE}} =
- \frac{1}{B} \sum_{i=1}^{B}
\log \frac{
\exp\!\left(\mathrm{sim}(z_i, z_i^+)/\tau\right)
}{
\sum_{j=1}^{B} \exp\!\left(\mathrm{sim}(z_i, z_j)/\tau\right)
}
\end{equation}
where $\mathrm{sim}(\cdot,\cdot)$ denotes cosine similarity and
$\tau$ is a temperature hyperparameter. After pre-training, the
encoder is fine-tuned end-to-end on the labelled downstream task.


% ═════════════════════════════════════════════════════════════════
\section{Experiments and Results}
\label{sec:experiments}
% ═════════════════════════════════════════════════════════════════

\subsection{Experimental Setup}

All models are implemented in PyTorch and trained on a CUDA-enabled
GPU. Signal processing uses SciPy (Butterworth filter, polyphase
resampling). Hyperparameter search uses Optuna with the TPE sampler
over 40 trials in a nested cross-validation framework, with Fold~1 as
the internal validation fold; remaining folds serve as holdout
validation. Edge inference runs on a Raspberry Pi CPU.

The base architecture uses $d_{\text{model}}{=}128$, $N{=}4$
cross-attention layers, $h{=}8$ attention heads, feed-forward
dimension $512$, dropout $0.1$, and sequence length $T{=}256$. The
edge-optimised variant uses $d_{\text{model}}{=}32$, $N{=}3$,
feed-forward dimension $256$, learning rate $2.91{\times}10^{-4}$,
weight decay $1.62{\times}10^{-4}$, and batch size $32$.

All results are reported as mean$\,\pm\,$standard deviation across
5-fold cross-validation. Evaluation metrics are accuracy, precision,
recall, and F1-score for each classification head, and combined
accuracy (average of the two head accuracies) as the primary
summary metric.

\subsection{Base Model Performance}

The base model trained with band-pass filtering and downsampling
achieved a mean accuracy of $93.31\%\pm0.22\%$ (precision
$93.81\%\pm0.31\%$, recall $93.31\%\pm0.22\%$, F1
$93.23\%\pm0.22\%$) for HC-PD classification and $92.34\%\pm0.08\%$
accuracy (precision $93.14\%\pm0.18\%$, recall $92.34\%\pm0.08\%$,
F1 $92.28\%\pm0.08\%$) for PD-DD classification across 5-fold
cross-validation, giving a combined accuracy of $92.83\%$. The higher
accuracy on the HC-PD partition confirms that differentiating PD from
DD is more challenging due to the commonality of motor patterns.
The t-SNE plot in Figure~\ref{fig:initial_tsne} shows a region of
overlap common to PD and DD samples that is absent in the HC-PD
embedding, visually confirming this difficulty.

% [TODO] ADD PER-FOLD BREAKDOWN TABLE HERE once per-fold numbers
%        are available. Structure:
% \begin{table}[htbp]
% \caption{Per-fold results, base model.}
% \begin{center}
% \begin{tabular}{cccc}
% \toprule
% Fold & HC-PD Acc & PD-DD Acc & Combined \\
% \midrule
% 1 & ? & ? & ? \\
% ... \\
% Mean$\pm$Std & $93.31\pm0.22$ & $92.34\pm0.08$ & $92.83\%$ \\
% \bottomrule
% \end{tabular}
% \end{center}
% \end{table}

\begin{figure}[!ht]
\centering
\begin{minipage}[b]{0.4\textwidth}
    \centering
    \includegraphics[width=\textwidth,height=0.2\textheight,keepaspectratio]{roc_hc_vs_pd.png}
    \centerline{\small (a) HC vs PD}
\end{minipage}
\hfill
\begin{minipage}[b]{0.4\textwidth}
    \centering
    \includegraphics[width=\textwidth,height=0.2\textheight,keepaspectratio]{roc_pd_vs_dd.png}
    \centerline{\small (b) PD vs DD}
\end{minipage}
\caption{ROC curves for the base model (Fold~1).}
\label{fig:initial_roc}
\end{figure}

\begin{figure}[!ht]
\centering
\begin{minipage}[b]{0.4\textwidth}
    \centering
    \includegraphics[width=\textwidth,height=0.2\textheight,keepaspectratio]{tsne_hc_vs_pd.png}
    \centerline{\small (a) HC vs PD}
\end{minipage}
\hfill
\begin{minipage}[b]{0.4\textwidth}
    \centering
    \includegraphics[width=\textwidth,height=0.2\textheight,keepaspectratio]{tsne_pd_vs_dd.png}
    \centerline{\small (b) PD vs DD}
\end{minipage}
\caption{t-SNE embeddings for the base model (Fold~1). The region of
         PD/DD overlap visible in~(b) is absent in the HC/PD
         embedding~(a).}
\label{fig:initial_tsne}
\end{figure}

\subsection{Band-Pass Filtering Ablation}
\label{sec:ablation_bp}

Band-pass filtering is fundamental in physiological signal processing.
Significant human motor activity such as tremor and FoG is concentrated
in the $0.1$--$20$~Hz band; signals above this range arise primarily
from sensor noise unrelated to motor symptoms.
Without band-pass filtering, the model achieved only
$85.64\%\pm2.14\%$ (HC-PD) and $85.39\%\pm1.88\%$ (PD-DD) accuracy
with substantially higher fold-to-fold variance, and the training
curves in Figure~\ref{fig:bandpass_loss} confirm severe overfitting.
With the filter applied, accuracy improved to $93.31\%\pm0.22\%$ and
$92.34\%\pm0.08\%$ respectively, with the fold-to-fold standard
deviation reduced by an order of magnitude, demonstrating that
removing out-of-band noise is essential for reliable learning.

\begin{figure}[!ht]
\centering
\begin{minipage}[b]{0.4\textwidth}
    \centering
    \includegraphics[width=\textwidth,height=0.2\textheight,keepaspectratio]{loss_no_bandpass.png}
    \centerline{\small (a) Without band-pass}
\end{minipage}
\hfill
\begin{minipage}[b]{0.4\textwidth}
    \centering
    \includegraphics[width=\textwidth,height=0.2\textheight,keepaspectratio]{loss.png}
    \centerline{\small (b) With band-pass}
\end{minipage}
\caption{Loss curves illustrating overfitting without band-pass
         filtering (Fold~1).}
\label{fig:bandpass_loss}
\end{figure}

\begin{figure}[!ht]
\centering
\begin{minipage}[b]{0.4\textwidth}
    \centering
    \includegraphics[width=\textwidth,height=0.2\textheight,keepaspectratio]{roc_3class.png}
    \centerline{\small (a) Three-class ROC}
\end{minipage}
\hfill
\begin{minipage}[b]{0.4\textwidth}
    \centering
    \includegraphics[width=\textwidth,height=0.2\textheight,keepaspectratio]{tsne_3class.png}
    \centerline{\small (b) Three-class t-SNE}
\end{minipage}
\caption{Three-class baseline ROC and t-SNE, showing significant
         PD/DD overlap.}
\label{fig:3class_eval}
\end{figure}

\subsection{Hierarchical vs.\ Three-Class Classification}

The three-class model achieved a mean accuracy of $88.35\%\pm0.83\%$
(precision $89.67\%\pm0.83\%$, recall $88.35\%\pm0.83\%$, F1
$88.32\%\pm0.88\%$). The learning objective of distinguishing healthy
samples (HC) from diseased samples (PD, DD) is qualitatively different
from distinguishing PD from DD. When optimised jointly, the model
fails to focus on the subtle difference between PD and DD samples,
resulting in lower performance on that boundary. The hierarchical
dual-head approach yields a $4.48$~pp improvement in combined
accuracy by decomposing the problem into two focused binary tasks.
The ROC curves in Figure~\ref{fig:3class_eval} and the t-SNE
embedding illustrate the model's difficulty in separating PD from DD
under the three-class objective.

\subsection{Hyperparameter Optimisation and Nested Cross-Validation}

Hyperparameter tuning was conducted to achieve regularisation and to
improve generalisation to unseen data. Hyperparameter search was
conducted using Optuna (TPE sampler, 40 trials, MedianPruner) on
Fold~1 as the internal validation fold within the nested
cross-validation framework; the remaining folds serve as holdout
validation to estimate generalisation performance without information
leakage from the search. This ensures that reported accuracy is a less
biased estimate of true out-of-sample performance.

\begin{table}[htbp]
\caption{Training vs.\ test accuracy at key epochs (Fold~1, Optuna
         optimised model).}
\label{tab:train_test_gap}
\begin{center}
\begin{tabular}{lccc}
\toprule
\textbf{Phase} & \textbf{Train Acc} & \textbf{Test Acc} & \textbf{Gap} \\
\midrule
Best checkpoint (Epoch~45) & 89.39\% & 91.82\% & $-$2.43\% \\
Final epoch (Epoch~80)     & 92.10\% & 88.78\% & $+$3.32\% \\
\bottomrule
\end{tabular}
\end{center}
\end{table}

The Optuna-optimised model (best checkpoint, Epoch~45) achieves a
test accuracy of $88.78\%$, approximately $4$~pp below the base model,
in exchange for a $4\times$ reduction in embedding dimension
($d{=}32$ vs.\ $d{=}128$). The negative train/test gap at Epoch~45
(test exceeds train) indicates good generalisation rather than
overfitting at that checkpoint. The optimised model shows
comparatively stronger performance on the PD-DD head, suggesting that
the smaller, more regularised architecture generalises the subtle
PD/DD boundary better.

\subsection{Edge Deployment on Raspberry Pi}
\label{sec:pi}

The optimised model was deployed on a Raspberry Pi CPU to evaluate
real-time inference. The inference pipeline downsamples incoming
$100$~Hz samples to $64$~Hz, applies a causal fifth-order Butterworth
band-pass filter sample-by-sample using second-order sections (SOS)
with persistent filter state, and accumulates $256$ samples in a
circular buffer. Inference is triggered every $128$ new samples
($50\%$ overlap, $2.0$~s step), producing a prediction at $0.5$~Hz.
The mean per-window inference time is $48.32$~ms at $0.02\%$ CPU
utilisation. This confirms feasibility for deployment on an edge
device such as a smartwatch for real-time monitoring.

\subsection{Self-Supervised Pre-training and Label Efficiency}

The self-supervised model fine-tuned on $100\%$ of labelled data
achieved a combined accuracy of $92.87\%$ (HC-PD: $93.35\%$,
PD-DD: $92.40\%$), competitive with the fully supervised base model.
A label-efficiency study evaluated fine-tuning across fractions of the
labelled set; results are summarised in Table~\ref{tab:ssl_label_eff}
and Figure~\ref{fig:label_eff}. With only $5\%$ of labelled data the
model reaches $80.49\%$ combined accuracy, and with $20\%$ it already
exceeds $91\%$. The rapid accuracy gain between $5\%$ and $20\%$,
followed by a gain of less than $2$~pp from $20\%$ to $100\%$,
demonstrates that the self-supervised pre-training has learned general
IMU representations that require very few labels to transfer
effectively to the downstream task.

\begin{table}[htbp]
\caption{SSL label-efficiency: fine-tuning accuracy vs.\ labelled
         fraction (combined accuracy across both heads).}
\label{tab:ssl_label_eff}
\begin{center}
\begin{tabular}{rccc}
\toprule
\textbf{Labels (\%)} & \textbf{HC-PD Acc} & \textbf{PD-DD Acc} &
\textbf{Combined} \\
\midrule
$5\%$   & $80.52\%$ & $80.46\%$ & $80.49\%$ \\
$10\%$  & $87.62\%$ & $86.12\%$ & $86.87\%$ \\
$20\%$  & $92.22\%$ & $89.82\%$ & $91.02\%$ \\
$50\%$  & $92.70\%$ & $91.69\%$ & $92.20\%$ \\
$70\%$  & $92.34\%$ & $92.10\%$ & $92.22\%$ \\
$100\%$ & $93.35\%$ & $92.40\%$ & $92.87\%$ \\
\bottomrule
\end{tabular}
\end{center}
\end{table}

\begin{figure}[!ht]
  \centering
  \includegraphics[width=0.45\textwidth,height=0.22\textheight,keepaspectratio]{accuracy_vs_labels.png}
  \caption{Combined accuracy vs.\ labelled fraction. Performance
           plateaus after $20\%$ labels, demonstrating effective
           representation transfer from self-supervised pre-training.}
  \label{fig:label_eff}
\end{figure}

\subsection{Comparison with Foundation Model (TimesFM + LoRA)}

For comparative analysis, we fine-tuned TimesFM~\citep{das2024}, a
foundation model pre-trained on large datasets for general-purpose
forecasting tasks. We adapted it for classification by adding a LoRA
(Low-Rank Adaptation) adapter~\citep{hu2022} with rank $r{=}6$, scaling
factor $\alpha{=}12$, dropout $0.15$, and regularisation $0.1$,
applied to the hidden dimension of $1{,}280$. LoRA is a
parameter-efficient fine-tuning technique that adapts large
pre-trained models to new tasks on smaller datasets by learning
low-rank update matrices. This achieved $89.70\%$ (HC-PD) and
$87.80\%$ (PD-DD) accuracy ($88.75\%$ combined), below the base
model. This suggests that a forecasting foundation model does not
transfer as effectively to this classification setting as a
task-specific architecture trained from scratch.

\subsection{Task-Wise Ablation (1D CNN and BiLSTM)}

We trained task-specific 1D CNN and BiLSTM models for comparison. The
1D CNN uses four convolutional blocks ($d{=}128$) with kernel sizes
$7, 5, 3, 3$, followed by batch normalisation and max pooling. The
BiLSTM uses two layers ($d{=}128$) with layer normalisation.

The BiLSTM showed higher accuracy for tasks involving sustained
postures with strong temporal structure: CrossArms ($92.39\%$),
HoldWeight ($93.87\%$), LiftHold ($93.92\%$), and StretchHold
($93.87\%$). The CNN performed comparatively better on tasks where
spatial patterns within a window are more informative than temporal
dependencies: Entrainment ($63.68\%$) and PointFinger ($65.86\%$).
Tasks such as Entrainment, PointFinger, TouchIndex, and TouchNose
consistently produced accuracy below $70\%$ for both models,
indicating that these tasks elicit motor patterns that are difficult
to discriminate at the window level for either architecture.
Figure~\ref{fig:taskwise} visualises the per-task comparison.

\begin{figure}[!ht]
  \centering
  \includegraphics[width=0.5\textwidth,height=0.22\textheight,keepaspectratio]{image6.png}
  \caption{Task-wise ablation: 1D CNN vs.\ BiLSTM combined accuracy
           per motor task.}
  \label{fig:taskwise}
\end{figure}

\subsection{Consolidated Results}

Table~\ref{tab:summary} consolidates 5-fold cross-validation results
across all experimental variants. The hierarchical dual-head base
model with band-pass filtering achieves the highest combined accuracy.
SSL pre-training matches the supervised baseline while requiring
significantly fewer labels.

\begin{table}[htbp]
\caption{Consolidated 5-fold cross-validation results
         (mean$\,\pm\,$std\,\%). Best combined accuracy in bold.}
\label{tab:summary}
\begin{center}
\begin{tabular}{lcccccc}
\toprule
\multirow{2}{*}{\textbf{Approach}} &
  \multicolumn{2}{c}{\textbf{HC-PD}} &
  \multicolumn{2}{c}{\textbf{PD-DD}} &
  \multirow{2}{*}{\textbf{Combined}} \\
\cmidrule(lr){2-3}\cmidrule(lr){4-5}
& \textbf{Acc} & \textbf{F1}
& \textbf{Acc} & \textbf{F1} & \\
\midrule
Base Model (w/ BP) &
  $93.31{\pm}0.22$ & $93.23{\pm}0.22$ &
  $92.34{\pm}0.08$ & $92.28{\pm}0.08$ & $\mathbf{92.83\%}$ \\
Base Model (w/o BP) &
  $85.64{\pm}2.14$ & --- &
  $85.39{\pm}1.88$ & --- & $85.52\%$ \\
Three-Class Baseline &
  \multicolumn{2}{c}{$88.35{\pm}0.83$} &
  \multicolumn{2}{c}{$88.32{\pm}0.88$} & $88.35\%$ \\
SSL Fine-tune (100\%) &
  $93.35$ & --- &
  $92.40$ & --- & $92.87\%$ \\
TimesFM LoRA ($r{=}6$) &
  $89.70$ & --- &
  $87.80$ & --- & $88.75\%$ \\
\midrule
Raspberry Pi (inference) &
  \multicolumn{4}{c}{---} & ${\approx}48.32$~ms/win \\
\bottomrule
\end{tabular}
\end{center}
\end{table}

% [TODO] ADD STATE-OF-THE-ART COMPARISON TABLE HERE once you have
%        identified prior published results on the PADS dataset.

\subsection{Cross-Attention Weight Analysis}

Figure~\ref{fig:attn} visualises the cross-attention weight maps to
provide insight into which temporal patterns the model attends to.
Attention bands become progressively sparser across the classes from
HC to PD and PD to DD. The HC class shows concentrated activity on
earlier timesteps and around timestep~250, while PD samples show more
distributed active bands, consistent with the irregular, recurrent
nature of tremor and bradykinesia. The similarity between the active
regions of PD and DD samples further confirms the difficulty of
separating Parkinson's from other neurodegenerative disorders at the
attention level.

\begin{figure}[!ht]
\centering
\begin{minipage}[b]{0.48\textwidth}
    \centering
    \includegraphics[width=\textwidth,height=0.2\textheight,keepaspectratio]{attention_comparison_hc_vs_pd.png}
    \centerline{\small (a) HC vs PD}
\end{minipage}
\hfill
\begin{minipage}[b]{0.48\textwidth}
    \centering
    \includegraphics[width=\textwidth,height=0.2\textheight,keepaspectratio]{attention_comparison_pd_vs_dd.png}
    \centerline{\small (b) PD vs DD}
\end{minipage}
\caption{Cross-attention weight maps (Fold~1). Sparser attention in
         PD and DD reflects distributed temporal motor symptom
         patterns.}
\label{fig:attn}
\end{figure}


% ═════════════════════════════════════════════════════════════════
\section{Discussion}
\label{sec:discussion}
% ═════════════════════════════════════════════════════════════════

The results demonstrate that the proposed Dual-Channel Cross-Attention
Transformer is well suited to the interdisciplinary challenge of PD
detection from wearable IMU data. The hierarchical binary
classification strategy consistently outperforms the three-class
baseline by decomposing a single hard objective into two focused tasks,
reflecting the clinical reality that HC/PD and PD/DD represent
qualitatively different discrimination problems. The $4.48$~pp
combined accuracy gain confirms that the model benefits from learning
each decision boundary independently.

The band-pass filtering ablation highlights a fundamental property of
medical time series data: the discriminative signal is concentrated in
a narrow, low-frequency, high-sensitivity band ($0.1$--$20$~Hz), and
out-of-band noise does not merely add difficulty---it actively prevents
the model from learning meaningful patterns. The substantially higher
fold-to-fold variance in the unfiltered condition ($\pm2.14\%$ vs.\
$\pm0.22\%$ for HC-PD) illustrates that noise inflates variance in
model estimates, which is precisely the regime where nested
cross-validation and careful signal processing matter most.

The self-supervised pre-training results reveal an important property
for clinical deployment: the model achieves over $91\%$ combined
accuracy with just $20\%$ of labelled data, and further gains from
additional labels are marginal. In practice, obtaining labelled
clinical data requires a physician or specialist to assess each
participant, which is a significant bottleneck. The SSL framework
demonstrates that general IMU representations can be learned without
any labels, substantially reducing the annotation burden for future
data collection in new clinical populations or settings.

The edge deployment on Raspberry Pi validates the feasibility of
real-time monitoring on constrained hardware. At $48.32$~ms per
inference window and $0.02\%$ CPU utilisation, the model is far within
the computational budget of a modern smartwatch processor, positioning
this framework as a viable component of a continuous disease monitoring
system in which PD progression can be tracked passively during daily
life.

The task-wise ablation reveals that no single baseline architecture is
uniformly superior across all motor tasks, and that some tasks
($<70\%$ for both models) remain genuinely difficult regardless of
architecture. This suggests that task-aware fusion strategies and
multi-task aggregation are important directions for improving
robustness.

\subsection{Limitations}

The primary limitation of this work is that all experiments are
conducted on a single dataset (PADS). While patient-level
cross-validation reduces optimistic bias within this dataset, external
validation on an independent cohort from a different clinical site
or smartwatch hardware is necessary to confirm generalisability.
Second, the self-supervised pre-training uses the same dataset as
fine-tuning; ideally, pre-training would leverage a larger pool of
unlabelled IMU data from diverse populations, which is the norm for
self-supervised learning in other domains. Third, clinical
validation---where model predictions are evaluated by a specialist
against ground-truth diagnosis over time---has not been performed and
remains essential before real-world deployment. Finally, the
Raspberry Pi results provide a proxy for smartwatch deployment but
not a direct evaluation; power consumption and latency on actual
smartwatch hardware should be measured in future work.


% ═════════════════════════════════════════════════════════════════
\section{Conclusion}
\label{sec:conclusion}
% ═════════════════════════════════════════════════════════════════

We proposed a Dual-Channel Cross-Attention Transformer framework for
detecting Parkinson's Disease from bilateral wrist-worn IMU data. By
processing left and right wrist signals through parallel encoders
connected by cross-attention, the model explicitly captures the motor
asymmetry that is a clinically established characteristic of PD. A
hierarchical binary classification strategy decomposes the
three-group problem into focused HC-PD and PD-DD tasks, yielding a
$4.48$~pp improvement over a three-class baseline. Patient-level
stratified 5-fold cross-validation provides less biased estimates of
true out-of-sample performance, and band-pass filtering is shown to
be an essential preprocessing step, reducing fold-to-fold variance by
an order of magnitude.

The self-supervised contrastive pre-training stage demonstrates that
competitive classification performance is achievable with only $20\%$
of labelled data, which is a critical finding for developing clinical
frameworks where expert annotation is scarce. The primary limitation
is reliance on a single dataset; external validation and direct
smartwatch hardware evaluation are required before clinical deployment.

Future work should investigate task-aware fusion strategies to improve
performance on difficult motor tasks such as Entrainment and
PointFinger. Expanding the framework to longitudinal data would enable
disease progression monitoring over time. The interdisciplinary nature
of this problem---spanning clinical neurology, signal processing, and
deep learning---means that clinical validation in partnership with
specialist assessment is an essential next step. Ultimately, a
validated continuous monitoring system of this kind could reduce the
burden on specialist clinics and enable earlier, more objective
intervention for millions of patients living with Parkinson's Disease.


% ─── References ──────────────────────────────────────────────────
% MLHC uses natbib. Use \citep{} for parenthetical and \citet{}
% for textual citations throughout.
% Replace the placeholder below with your actual .bib file.
\bibliography{references}

% ─── Appendix (optional, does not count toward page limit) ───────
% \appendix
% \section{Additional Results}

\end{document}